
\begingroup
\fontsize{12}{12}\selectfont

\renewcommand{\arraystretch}{1.3}
\setlength{\tabcolsep}{5pt}
\setlength{\fboxsep}{3pt}
\setlength{\fboxrule}{0.5pt}

\noindent\fbox{%
	\begin{minipage}{\dimexpr\linewidth-2\fboxrule\relax}
		
		\rowcolors{1}{gray!10}{white}
		
		\begin{tabularx}{\linewidth}{
				>{\raggedright\arraybackslash}p{0.38\linewidth}
				>{\raggedright\arraybackslash}X
			}
			
			\texttt{00\_FinalPoster}$^*$ &
			\,\tabitem Die allerletzte Version Ihres Posters als pdf mit Band, Jahr
			und Seitenzahlen auf dem File-Server unter 'Publikationen'
			\par\smallskip
			\,\tabitem EndNote-Eintrag für die IWT-WiKo mit DOI-Nummer, und ein EndNote
			*.ris ASCII-Export, der den FAIR-Prinzipien entspricht \\
			
			\texttt{01\_Poster+Abstract} &
			\,\tabitem das *.pdf und *.pptx or *.tex \\
			
			\texttt{02\_Figures}$^{**}$ &
			\,\tabitem ein Unterverzeichnis für jede Abbildung (figure\_01, ...) mit
			Datenpunkten als*.csv or *.txt, oder die erzeugende Datei/ das
			erzeugende Programm für die Abbildungspunkte (Matlab, Python,
			Origin,...); die Ergebnisse stammen aus Daten, die im
			Verzeichnis \texttt{04\_PrimaryData} gespeichert oder daraus
			extrahiert wurden \\
			
			\texttt{03\_Tables}$^{**}$ &
			\,\tabitem ein Unterverzeichnis für jede Tabelle (Tabelle\_01, ...) in dem
			die Daten in Form einer csv-Datei gespeichert sind \\
			
			\texttt{04\_PrimaryData} &
			\,\tabitem gemessene oder simulierte Daten, jedes Unterverzeichnis mit
			einer “readme.json” Datei zur Beschreibung der Daten
			\par\smallskip
			\,\tabitem ELN pdf mit QR-code (siehe unten) \\
			
			\texttt{05\_ProgramSources} &
			\,\tabitem e.g. OpenFOAM-Solver, Python/Matlab code, etc.
			\par\smallskip
			\,\tabitem Fallbeispiele (Netze) for OpenFOAM, Fluent, ABAQUS etc.
			\par\smallskip
			\,\tabitem “readme.json” Datei mit Metadaten \\
			
			\texttt{06\_References} &
			\,\tabitem die Referenzen *.bib oder*.enl zusammen mit *.ris
			\par\smallskip
			\,\tabitem Zitierte Arbeiten/Bücher (wenn Veröffentlichungsrechte
			erlaubt sind) \\
			
			\texttt{07\_Supplement} &
			\,\tabitem Ergänzende Dateien \\
			
			\texttt{08\_Misc} &
			\,\tabitem e.g. das Konferenzprogramm \\
			
		\end{tabularx}
		
	\end{minipage}%
}

\endgroup
